\section*{Capítulo II, Problema X}

\textbf{Problema:}

10. Sea $P = (1,2,3,4)$ y $Q = (4,3,2,1)$. Sea $A = (1,1,1,1)$. Sea $L$ la recta que pasa por $P$ y es paralela a $A$.

\begin{itemize}
    \item[(a)] Dado un punto $X$ en la recta $L$, calcula la distancia entre $Q$ y $X$ como función del parámetro $t$.
    \item[(b)] Demuestra que existe precisamente un punto $X_0$ en la recta para el cual esta distancia alcanza un mínimo, y que este mínimo es $2\sqrt{5}$.
    \item[(c)] Demuestra que $X_0 - Q$ es perpendicular a la recta.
\end{itemize}

\textbf{Solución:}

La recta $L$ que pasa por $P$ y es paralela a $A$ puede escribirse como
\[
    X(t) = P + tA = (1+t, 2+t, 3+t, 4+t).
\]

\begin{itemize}
    \item[(a)] La distancia entre $Q$ y $X(t)$ es
    \[
        d(t) = \|X(t) - Q\|.
    \]
    Como
    \[
        X(t) - Q = (t-3, t-1, t+1, t+3),
    \]
    se tiene
    \begin{align*}
        d(t)^2 &= (t-3)^2 + (t-1)^2 + (t+1)^2 + (t+3)^2 \\
               &= 4t^2 + 20.
    \end{align*}
    Por tanto,
    \[
        d(t) = \sqrt{4t^2 + 20} = 2\sqrt{t^2 + 5}.
    \]

    \item[(b)] Como $d(t) = 2\sqrt{t^2+5}$, la distancia es mínima exactamente cuando $t^2$ es mínima. Esto ocurre únicamente en $t=0$.

    Luego el único punto de mínima distancia es
    \[
        X_0 = X(0) = P = (1,2,3,4).
    \]
    La distancia mínima vale
    \[
        d(0) = 2\sqrt{5}.
    \]

    \item[(c)] La dirección de la recta es el vector
    \[
        A = (1,1,1,1).
    \]
    Además,
    \[
        X_0 - Q = (1,2,3,4) - (4,3,2,1) = (-3,-1,1,3).
    \]
    Calculamos su producto punto con el vector dirección de la recta:
    \[
        (-3,-1,1,3) \cdot (1,1,1,1) = -3 - 1 + 1 + 3 = 0.
    \]
    Por lo tanto, $X_0 - Q$ es perpendicular a la recta $L$.
\end{itemize}